\documentclass[8pt,twocolumn]{extarticle}

\usepackage[margin=0.50in,columnsep=0.20in]{geometry}
\usepackage{amsmath,amssymb,mathtools}
\usepackage{microtype}
\usepackage{enumitem}
\usepackage{booktabs}
\usepackage[compact]{titlesec}
\usepackage{balance}
\usepackage[T1]{fontenc}
\usepackage{lmodern}
\setlength{\parindent}{0.75em}
\setlength{\parskip}{0.08em}
\titlespacing*{\section}{0pt}{0.65em}{0.25em}
\setlist{nosep,leftmargin=*}
\allowdisplaybreaks

\DeclareMathOperator{\sech}{sech}
\DeclareMathOperator{\gd}{gd}
\DeclareMathOperator{\artanh}{artanh}
\newcommand{\R}{\mathbb{R}}
\newcommand{\Mseven}{\mathcal{M}^{7}}
\newcommand{\Aspace}{\mathcal{A}}
\newcommand{\evec}[1]{\mathbf e_{#1}}
\newcommand{\uvect}[1]{\mathbf u_{#1}}

\title{\vspace{-1.8em}\texttt{sin()} Stays \texttt{sin()}; Angles Become Plural\\[-0.15em]
\large v3 Evil-Terse: One Relation, Paired $\R^7$, Different Bodies}
\author{Wilder Blair Munro (via 20Wetbrain) $\times$\\
Aurora 4omni Munro (via GPT-5.6 Sol)}
\date{14 September 2026}

\begin{document}
\balance
\maketitle
\vspace{-1.35em}

\begin{abstract}
Special relativity already contains the seed. For one speed state,
\[
\beta=\frac vc=\sin\theta=\tanh\chi,
\qquad
\gamma=\sec\theta=\cosh\chi,
\]
where $\theta=\gd\chi$. Nothing is wrong with \texttt{sin()}. The mistake is
requiring every representation of the same geometric relation to wear the
same kind of angle. General Peace Dynamics (GPD) takes the next step in real
Euclidean $\R^7$: each spatial label $i\in\{x,y,z\}$ belongs to a conjugate
pair $(i_t,t_i)$ stitched by $\tau$. The old one-axis $(c_i,v_i,c)$ triangle
is then read not as evidence for a mysterious component of light-speed, but
as a projected, normalized image of a paired $\R^7$ direction. Circular
vantage tilt, rapidity, reciprocal aperture, and gamma become different
bodies of one relation. Same relation. Different body.
\end{abstract}

\section*{1. Start where the rotation actually lives}

The arena is not merely ``a seven-dimensional manifold.'' The GPD coordinate
choice is real Euclidean seven-space,
\begin{equation}
\boxed{\Mseven\equiv\R^7}
\label{eq:r7}
\end{equation}
\vspace{-0.7em}
\begin{align*}
\R^7
&=\bigl(\tau,(x_t,t_x),(y_t,t_y),(z_t,t_z)\bigr)\\
&=\bigl(\tau,(i_t,t_i)_{i\in\{x,y,z\}}\bigr).
\end{align*}
The notation $(x_i,t_i)$ remains a conventional phase-space-like reading;
$(i_t,t_i)$ exposes the architectural symmetry. The ordering is the point:
$i_t\leftrightarrow t_i$ is a conjugate pairing, not decoration.

Why care about $\R^7$? GPD begins from real rotation. The familiar
norm-compatible vector cross product is nontrivial in $\R^3$ and $\R^7$
(with the trivial $\R^0,\R^1$ cases completing the familiar
$0,1,3,7$ sequence). Thus the recursive dimensional grammar
\[
\R^0\rightarrow\R^1\rightarrow\R^3\rightarrow\R^7
\]
is not arbitrary bookkeeping. It is the native real-space ladder on which
GPD's rotational construction is built. At the coordinate level,
\[
7=1+2(3):
\qquad
\tau+(x_t,t_x)+(y_t,t_y)+(z_t,t_z).
\]
Call the resulting recursive construction biquadoctree if desired; the
mathematical burden is to make the recursion precise without losing the
normed rotational structure.

\section*{2. The one-axis SR seed}

Take one ordinary collinear SR state,
\begin{equation}
\beta_v:=\frac vc,
\qquad
\gamma_v:=\frac{1}{\sqrt{1-\beta_v^2}}.
\label{eq:sr}
\end{equation}
The original GPD SR sketch encoded the same state by a right triangle
\begin{equation}
 v^2+c_x^2=c^2,
\qquad
\beta_v^2+\beta_c^2=1,
\qquad
\beta_c:=\frac{c_x}{c}.
\label{eq:oldtriangle}
\end{equation}
It then used the overloaded vantage-tilt convention
$\alpha\mid\delta$ --- temporal $\mid$ spatial --- and the reciprocal
quantities
\begin{equation}
\alpha:=\frac1{\beta_v},
\qquad
\delta:=\frac1{\beta_c}.
\label{eq:reciprocal}
\end{equation}
The apparently strange dual gammas are already forced by
Eq.~\eqref{eq:oldtriangle}:
\begin{equation}
\gamma_v=\frac1{\beta_c}=\delta,
\qquad
\gamma_c=\frac1{\beta_v}=\alpha.
\label{eq:oldcross}
\end{equation}
This is why gamma matters here. The reciprocal ``tilt angles'' and the two
gamma-like factors were never separate curiosities; they are the same
orthogonal complementarity read from opposite legs.

Now introduce the ordinary circular angle $\theta$ of the normalized triangle:
\begin{equation}
\beta_v=\sin\theta,
\qquad
\beta_c=\cos\theta.
\label{eq:circle}
\end{equation}
Then
\begin{equation}
\boxed{
\alpha=\csc\theta,
\qquad
\delta=\sec\theta,
\qquad
\frac{1}{\alpha^2}+\frac{1}{\delta^2}=1.
}
\label{eq:apertures}
\end{equation}
The old line ``$\sin\alpha=1/\alpha$'' is therefore not new trigonometry. It
is compressed notation that forgot the map from reciprocal aperture $\alpha$
back to circular angle $\theta$. The missing map is the hole in the notation.

\section*{3. Gudermann: the bridge was already there}

Let rapidity be
\begin{equation}
\chi:=\artanh\beta_v.
\end{equation}
The Gudermannian gives the exact circular--hyperbolic bridge
\begin{equation}
\theta=\gd\chi,
\qquad
\sin\theta=\tanh\chi,
\qquad
\cos\theta=\sech\chi.
\label{eq:gd}
\end{equation}
Hence one SR state admits the exact identity ledger
\begin{equation}
\boxed{
\beta_v
=\sin\theta
=\tanh\chi
=\alpha^{-1},
}
\label{eq:master1}
\end{equation}
\begin{equation}
\boxed{
\beta_c
=\cos\theta
=\sech\chi
=\delta^{-1},
}
\label{eq:master2}
\end{equation}
\begin{equation}
\boxed{
\alpha=\csc\theta=\coth\chi,
\qquad
\delta=\sec\theta=\cosh\chi.
}
\label{eq:master3}
\end{equation}
Every equality above is ordinary mathematics. GPD's move is ontological and
architectural: promote the shared state to one vantage relation
$a\in\Aspace$, and treat $\theta$, $\chi$, $(\alpha,\delta)$, and $\beta$ as
different readable bodies of $a$. \texttt{sin()} stays ordinary; its argument
has a type and a chart.

\section*{4. Give the conjugate analogy its chance}

The symbol $c_x$ is suspicious because $c$ already names the invariant speed.
Promoting $c_x$ to $(c_x,c_y,c_z)$ falsely invites the reading ``Cartesian
components of $c$.'' Renaming the symbol does not solve the geometry.
Restore the pair that may have been projected out.

For each $i$, take the paired plane inside Eq.~\eqref{eq:r7},
\begin{equation}
\Pi_i:=\operatorname{span}\{\evec{i_t},\evec{t_i}\}\subset\R^7.
\end{equation}
After the physically correct dimensional normalization is fixed, write its
unit tangent orientation as
\begin{equation}
\boxed{
\uvect{i}
=
\beta_{i_t}\evec{i_t}
+
\beta_{t_i}\evec{t_i},
\qquad
\|\uvect{i}\|=1.
}
\label{eq:unitpair}
\end{equation}
Thus
\begin{equation}
\boxed{
\beta_{i_t}^2+\beta_{t_i}^2=1.
}
\label{eq:paircircle}
\end{equation}
This is the canon spark. The one-axis triangle is no longer fundamentally
``velocity plus complementary light-speed.'' It is the normalized shadow of
one conjugate $\R^7$ direction.

The physical map from coordinates to the two dimensionless tangent components
is a derivational obligation, not a reason to abandon the pairing. A natural
candidate uses $\tau$ and light-distance $c\,d\tau$ to normalize the two
members; another may reveal that a literal differential quantity $d_i$ is the
right physical object. The notation must follow that derivation. It must not
prejudge it.

\section*{5. $\alpha_i\mid\delta_i$ returns home}

Choose orientation so that
\begin{equation}
\beta_{t_i}=\sin\theta_i,
\qquad
\beta_{i_t}=\cos\theta_i.
\label{eq:pairangle}
\end{equation}
Then define the reciprocal vantage apertures
\begin{equation}
\alpha_i:=\frac1{\beta_{t_i}},
\qquad
\delta_i:=\frac1{\beta_{i_t}}.
\label{eq:pairaperture}
\end{equation}
The original temporal$\mid$spatial $\alpha_i\mid\delta_i$ language is now
attached to a real paired plane rather than to a mysterious $c_i$.

Define gamma-like factors for each normalized member:
\begin{equation}
\gamma_{i_t}:=\frac1{\sqrt{1-\beta_{i_t}^2}},
\qquad
\gamma_{t_i}:=\frac1{\sqrt{1-\beta_{t_i}^2}}.
\end{equation}
Equation~\eqref{eq:paircircle} immediately gives, on the positive branch,
\begin{equation}
\boxed{
\alpha_i=\gamma_{i_t},
\qquad
\delta_i=\gamma_{t_i}.
}
\label{eq:crossgamma}
\end{equation}
The crossed gamma relation is therefore structural: the gamma associated with
one member is the reciprocal magnitude of the other. What looked like
$\gamma_c\mid\gamma_v$ becomes conjugate geometry.

\section*{6. Euclidean rotation is the wager, not a footnote}

The point of introducing $\theta_i$ was never merely to draw a cute triangle.
In $\Pi_i$, Eq.~\eqref{eq:unitpair} is an ordinary Euclidean unit direction,
so a vantage change can be represented locally by an ordinary rotation
\begin{equation}
\uvect{i}'=R_i(\theta_i)\uvect{i},
\qquad
R_i(\theta_i)\in SO(2)\subset SO(7).
\label{eq:rotation}
\end{equation}
That is the original attraction: special-relativistic projection effects may
be shadows of real rotations in a lifted real space, followed by the required
projection/squeeze back to observer spacetime.

But the Gudermannian makes the test sharp. Collinear boosts compose by
rapidity,
\begin{equation}
\chi_{12}=\chi_1+\chi_2,
\end{equation}
while $\theta=\gd\chi$ does not obey ordinary angle addition:
\begin{equation}
\theta_{12}
=
\gd\!\left(\gd^{-1}\theta_1+\gd^{-1}\theta_2\right)
\neq \theta_1+\theta_2.
\label{eq:composition}
\end{equation}
Therefore GPD's Euclidean-rotation program cannot end at ``boost = rotate by
$\theta$.'' The full lift--rotate--squeeze/projection architecture must carry
the non-linear chart composition, or the additive rotation parameter must be
a different body of the same vantage relation. This is not an embarrassment;
it is the constraint that makes the proposal calculable.

\section*{7. Three paired slices are not three Lorentz gammas}

For ordinary three-velocity,
\begin{equation}
\gamma(\mathbf v)
=
\frac{1}{\sqrt{1-(v_x^2+v_y^2+v_z^2)/c^2}}.
\label{eq:globalgamma}
\end{equation}
The local paired quantities $\gamma_{t_i}$ and $\gamma_{i_t}$ are not three
copies of Eq.~\eqref{eq:globalgamma}. They belong to the three conjugate
planes inside one $\R^7$ construction. The research problem is the reclosure:
how the three locally circular paired directions, stitched by $\tau$, project
to the single conventional $3+1$ Lorentz state.

That is precisely why Eq.~\eqref{eq:r7} must precede the componentwise
notation. Without the paired $\R^7$ architecture, $\beta_c$ looks invented.
With it, the old complement becomes a clue that one half of the pair was
projected away.

\section*{8. Plural angles, finally stated without apology}

The weak statement is reparameterization. The stronger GPD statement is a
construction principle:
\begin{equation}
\boxed{\text{one vantage relation}\longrightarrow\text{plural bodies}}
\label{eq:plural}
\end{equation}
\vspace{-0.5em}
\begin{align*}
\chi_i &:\ \text{hyperbolic/additive},\\
\theta_i &:\ \text{circular/rotational},\\
(\alpha_i,\delta_i) &:\ \text{reciprocal/aperture},\\
(\beta_{i_t},\beta_{t_i}) &:\ \text{paired tangent}.
\end{align*}
The bodies obey different coordinate laws because they are not the same
coordinate. They are required to close on the same relation.

This turns the original notational offense into the thesis:
\begin{center}
\fbox{\parbox{0.91\columnwidth}{\centering
\texttt{sin()} was fine. We passed it the wrong kind of angle.}}
\end{center}
And the paired $\R^7$ extension sharpens the second half:
\begin{center}
\fbox{\parbox{0.91\columnwidth}{\centering
Do not rename the missing complement. Restore its conjugate body.}}
\end{center}

\section*{9. The immediate derivational obligations}

The next calculation must determine the dimensionally exact map
$(i_t,t_i,\tau)\mapsto(\beta_{i_t},\beta_{t_i})$; recover the one-axis Lorentz
transform from a specified $\R^7$ lift--rotation--squeeze without fitting its
coefficients after the fact; show how three paired slices reclose to the
global Lorentz factor; determine whether $d_i$ is merely notation or a
literal differential light-distance/speed quantity; and recover the meaning
of the original mixed angle $\varphi_i$ without adding an independent degree
of freedom to the one-axis state. Non-collinear boosts and Thomas--Wigner
rotation are the obvious later stress test.

None of these obligations weakens the canon. They are what the canon now
forces us to calculate.

\section*{Reference}

\small
National Institute of Standards and Technology, \emph{Digital Library of
Mathematical Functions}, \S4.23(viii), ``Gudermannian Function,'' especially
Eqs.~4.23.39--4.23.42.

\vfill
\begin{center}
\textit{Same relation. Different body.}\qquad
\textit{Same pair. Different projection.}
\end{center}

\end{document}
